\section{Rigidity, reconstruction, and the converse}
\label{k:section}

Assume equality on the strict area branch. Set
\[
 R_0=\sqrt{\frac{A}{4\pi}},\qquad M=m,\qquad
 a=-\frac JM,\qquad
 b=\sqrt{M^2-a^2-Q^2},\qquad \kappa=\frac{b}{R_0^2}>0.
\]
The parameter identities in Lemma~\ref{a:model-parameters} give
\((M+b)^2+a^2=R_0^2\). Until the reconstruction below, \(g,k\)
denote the polarized data, \(h\) their meridional metric, and
\(X=|\eta|\). We use the lapse \(N\), meridional shift \(Y\), and
stationary map \(\Phi\) supplied by Proposition~\ref{j:adjoint}.
Thus
\[
 H_3=-N^2\dd t_0^2+h_{ij}(\dd x^i+Y^i\dd t_0)
                  (\dd x^j+Y^j\dd t_0),\qquad \gamma=X^2H_3
\]
satisfies the Einstein--wave-map equations \eqref{a:wave}.
In particular, the argument starts with the full quotient equations
and their initial momenta.

\subsection{Strict timelikeness of the quotient Killing field}

Define
\begin{equation}\label{k:norm}
 \mu_*=N^2-|Y|_h^2,\qquad \alpha_Y=h(Y,\cdot).
\end{equation}
Here \(\mu_*\) is a Killing norm, unrelated to a constraint energy
density. Ordinary axial parity extends \(h,N,Y,\mu_*\) smoothly across
the axes to the doubled meridian. This is a planar annulus with a
smooth compact inner boundary; the transverse component of \(Y\) is
odd at an axis. The original open meridian is simply connected.

\begin{lemma}\label{k:timelike}
The function \(\mu_*\) is positive throughout the interior, including
on the axes. There is a globally defined function \(f\) on the open
meridian such that
\begin{equation}\label{k:static}
 \dd f=-\frac{\alpha_Y}{\mu_*},\qquad
 H_3=-\lambda^2\dd t^2+h_+,\qquad
 \lambda=\sqrt{\mu_*},\qquad
 h_+=h+\frac{\alpha_Y^2}{\mu_*},\qquad t=t_0+f.
\end{equation}
The original horizontal slice is the graph \(t=f\).
Moreover \(\rho=X\lambda\) is \(h_+\)-harmonic off the axes.
\end{lemma}

\begin{proof}
Let \(\xi=\partial_{t_0}\). Stationarity of \(\Phi\) and
\(\Ric_\gamma=2\Phi^*G\) imply \(\Ric_\gamma(\xi,\cdot)=0\).
The three-dimensional Killing twist
\[
 \mathcal W=*_\gamma
     \bigl(\xi^\flat_\gamma\wedge\dd\xi^\flat_\gamma\bigr)
\]
is a scalar. Differentiating with the parallel volume tensor, using
\(\nabla_i\xi_j=-\nabla_j\xi_i\) and the Killing second-derivative
identity, gives
\(\dd\mathcal W=\pm2*_\gamma
(\xi^\flat_\gamma\wedge\Ric_\gamma(\xi,\cdot))=0\).
Under \(\gamma=X^2H_3\), the covector wedge acquires the factor \(X^4\)
and the Hodge star on three-forms the factor \(X^{-3}\). Consequently
\(\mathcal W\) is \(X\) times the corresponding twist for \(H_3\).
The latter is bounded near any ordinary interior axis point:
\(H_3\) is smooth and nondegenerate there because \(N>0\).
Since \(X\) vanishes there, the constant \(\mathcal W\) is zero.
Substituting \(\xi^\flat_{H_3}=-\mu_*\dd t_0+\alpha_Y\) into
Frobenius gives
\begin{equation}\label{k:frobenius}
 \mu_*\,\dd\alpha_Y=\dd\mu_*\wedge\alpha_Y .
\end{equation}
This identity extends smoothly to the doubled meridian.

Where \(\mu_*>0\), Equation~\eqref{k:frobenius} makes
\(-\alpha_Y/\mu_*\) closed and gives the local static form
\eqref{k:static}. In these coordinates
\(\gamma=-\rho^2\dd t^2+X^2h_+\).
The \(tt\) Ricci equation is
\(\rho\Delta_{X^2h_+}\rho=0\); conformal invariance of the
two-dimensional scalar Laplacian then gives
\(\Delta_{h_+}\rho=0\).

We next exclude an interior null set. First there is a positive
collar of the inner boundary. The boundary alternatives from
Proposition~\ref{j:adjoint} are exhaustive. If \(N>0\) on \(S\), then
\(Y=N\nu\), \(\operatorname{sym}\nabla Y=-Nk\), and
\eqref{j:boundary} give
\[
 \partial_\nu\mu_*
 =2N\bigl(\partial_\nu N+Nk(\nu,\nu)\bigr)
 =2N\kappa>0.
\]
If \(N=0\) on \(S\), then \(Y=0\) there. Its tangential derivatives
vanish, and the symmetric-gradient equation makes its remaining
first derivatives vanish as well. At inward boundary distance \(s\),
\[
 Y=O(s^2),\qquad N=\kappa s+O(s^2),\qquad
 \mu_*=\kappa^2s^2+O(s^3)>0\quad(s>0\text{ small}).
\]
These collars can be chosen equivariantly and are uniform through
the poles by compactness. The end also has \(\mu_*>0\), since
\(N\to1\) and \(Y\to0\).

Suppose that the compact interior set
\(\mathcal Z=\{\mu_*=0\}\) on the double is nonempty.
Proposition~\ref{j:adjoint} gives \(\mu_*\ge0\) everywhere and
\(N>0\) in the interior, so \(Y\ne0\) on \(\mathcal Z\).
Let \(Z\) be the \(h\)-orthogonal quarter-turn of \(Y\), using the
orientation of the double. Evaluation of \eqref{k:frobenius} on
\((Y,Z)\) yields, on a neighborhood of \(\mathcal Z\),
\begin{equation}\label{k:null-flow}
 Z(\mu_*)=-\frac{\dd\alpha_Y(Y,Z)}{|Y|_h^2}\,\mu_* .
\end{equation}
The coefficient is smooth. Uniqueness for this scalar ordinary
differential equation shows that the flow of \(Z\) preserves
\(\mathcal Z\). A trajectory starting there exists for all time,
because it stays in a compact set on which \(Z\) is smooth.
Its omega-limit set has no equilibrium. Extend \(Z\) smoothly in
the plane around this compact set and apply the planar
Poincar\'e--Bendixson theorem: \(\mathcal Z\) contains a periodic
orbit, a simple closed curve \(C\).

Every component of \(\{\mu_*>0\}\) on the bounded side of \(C\),
intersected with either open regular half-meridian, would carry a
positive harmonic function \(\rho=X\sqrt{\mu_*}\).
This function extends continuously to the compact closure and
vanishes on its entire frontier: the frontier lies in an axis,
the inner boundary, or \(\mathcal Z\). A positive maximum would
therefore occur at an interior point, where \(h_+\) is smooth and
elliptic. The local strong maximum principle gives a contradiction.
No ellipticity at the frontier is needed in this argument.
Intersections of \(C\) with a reflected axis cause no additional
frontier: first restrict to either open half and then take a
component; its relative frontier is contained in the listed sets,
with \(C\) itself already contained in \(\mathcal Z\).

If \(C\) surrounds the inner obstacle, the positive collar supplies
such a component. Otherwise \(C\) bounds a disk disjoint from the
obstacle. The same maximum argument, together with \(\mu_*\ge0\),
forces \(\mu_*=0\) on that entire disk; a positive point on an axis
would have positive off-axis neighbors. Hence \(Y\), and therefore
\(Z\), is nonvanishing on the closed disk. Its restriction to \(C\)
is a tangent field to a simple periodic orbit and has winding
number one, whereas a nonvanishing field extending over a disk has
winding number zero. This last contradiction proves
\(\mu_*>0\) throughout the interior. Simple connectivity now gives
the global primitive \(f\) in \eqref{k:static}.
\end{proof}

\subsection{The static boundary and its conformal scale}

\begin{lemma}\label{k:static-boundary}
The static meridian has a smooth compact inner boundary, with
\[
 h_+|_{TS}=h|_{TS},\qquad \lambda=0,\qquad
 |\dd\lambda|_{h_+}=\kappa .
\]
If \(N|_S=0\), then its boundary coordinates and \(f\) are smooth in
the original one-sided coordinates. If \(N|_S>0\), then
\(\lambda=\sqrt{\mu_*}\) is a smooth static boundary coordinate,
the static metric has an isometric reflection across
\(\lambda=0\), and
\begin{equation}\label{k:log-time}
 f=-\frac{\log\mu_*}{2\kappa}+f_0
\end{equation}
with \(f_0\) smooth in the original coordinates up to \(S\).
These statements are compatible with ordinary axis regularity at
the two poles.
\end{lemma}

\begin{proof}
In the zero-lapse case, the preceding Taylor expansions give
\(\mu_*=s^2(\kappa^2+O(s))\). Taylor division on \(s\ge0\) shows that
\(\lambda=s(\kappa+O(s))\), \(\alpha_Y/\mu_*\), and
\(\alpha_Y^2/\mu_*\) are smooth, the last vanishing to order two.
Thus \(h_+\) extends smoothly with the stated boundary value and
normal derivative, and \(\dd f\) extends smoothly. Its primitive
also extends, after fixing an additive constant. This argument
only requires one-sided smoothness.

In the positive-lapse case, \(\mu_*\) itself is a boundary defining
coordinate. At \(S\),
\(\dd\mu_*=2N\kappa\,\nu^\flat\) and
\(\alpha_Y=N\nu^\flat\), so in a smooth collar \((\mu_*,z)\) we can
write
\[
 \alpha_Y=a_0\,\dd\mu_*+\mu_*\eta_0,\qquad
 a_0=\frac1{2\kappa}+\mu_*a_1 ,
\]
where \(a_1\) and the tangential one-form \(\eta_0\) are smooth.
Pull back by \(\mu_*=\lambda^2\). Then
\begin{equation}\label{k:reflected-base}
 h_+=h+\bigl(2a_0\,\dd\lambda+\lambda\eta_0\bigr)^2 .
\end{equation}
The pulled-back \(h\) has even tangential and normal coefficients
and odd mixed coefficients under \(\lambda\mapsto-\lambda\).
The one-form in parentheses changes sign under this reflection.
The metric is consequently invariant, with vanishing mixed
coefficients and normal coefficient \(4a_0^2=\kappa^{-2}\) at the
boundary. Its tangential block is the positive boundary block of
\(h\), proving nonsingularity and the asserted derivative of
\(\lambda\). Moreover
\[
 \dd f=-\frac{\dd\mu_*}{2\kappa\mu_*}
                   -a_1\dd\mu_*-\eta_0.
\]
The smooth remainder is closed, because \(\dd f\) is closed.
Integration in the collar gives \eqref{k:log-time}.

Choose all collars equivariantly under meridional axis reflection.
The preceding divisions then preserve even and odd polar parity
also at the poles. On an ordinary axis away from its endpoints
the added term in \(h_+\) has zero transverse radial component.
Thus the usual no-cone condition is retained:
\(|\dd X|_{h_+}=1\) on such an axis.

We make the joint corner regularity explicit. Near a pole in
the positive-lapse case choose equivariant coordinates
\((\mu_*,y)\), with \(y\) a signed transverse axis coordinate. The signed
extension of \(X\) is
\[
 X=yA_X(\mu_*,y^2),\qquad A_X>0.
\]
The coefficients \(h_{\mu_*y}\) and \(\eta_y\), where
\(\eta_0=\eta_y\dd y\), are odd in \(y\), while
\(h_{yy}(\mu_*,0)=A_X(\mu_*,0)^2\) by the original no-cone
condition. In \((\lambda,y)\), Equation~\eqref{k:reflected-base}
has components
\[
 \begin{split}
 (h_+)_{\lambda\lambda}&=4\lambda^2h_{\mu_*\mu_*}+4a_0^2,\\
 (h_+)_{\lambda y}&=2\lambda h_{\mu_*y}+2a_0\lambda\eta_y,\\
 (h_+)_{yy}&=h_{yy}+\lambda^2\eta_y^2.
 \end{split}
\]
They extend smoothly under both commuting reflections
\(\lambda\mapsto-\lambda\), \(y\mapsto-y\).
At the corner the metric is diagonal and positive, with normal
entry \(\kappa^{-2}\), and along \(y=0\) its \(yy\) entry is
\(A_X(\lambda^2,0)^2\). Thus no-cone regularity persists jointly
up to the pole. The zero-lapse case uses the original smooth
boundary coordinate \(s\) and the same \(y\)-reflection, without
requiring an \(s\)-reflection.
\end{proof}

\begin{lemma}\label{k:rod-scale}
There are global static isothermal coordinates
\((\sigma,\vartheta)\in[0,\infty)\times[0,\pi]\), with the inner
boundary \(\sigma=0\) and the two oriented axes
\(\vartheta=0,\pi\), in which
\begin{equation}\label{k:rho}
 h_+=P_+(\dd\sigma^2+\dd\vartheta^2),\qquad
 \rho=b\sinh\sigma\sin\vartheta .
\end{equation}
The boundary area density is
\begin{equation}\label{k:area-density}
 X\sqrt{P_+}=\frac b\kappa\sin\vartheta
                 =R_0^2\sin\vartheta\qquad(\sigma=0).
\end{equation}
In particular, the scale \(b\) is determined by the original area
and the adjoint boundary condition.
\end{lemma}

\begin{proof}
Reflect the static meridian across the axes. Its conformal filling
at infinity is a disk with a smooth boundary. Uniformize this
disk, sending infinity to its center and the axis reflection to
complex conjugation; inversion gives an exterior circular
coordinate, and its logarithm gives the asserted strip.
Smooth isothermal and boundary conformal regularity give
nonsingular charts through the ordinary axes and up to the inner
boundary, including their intersections. This uses the smooth
boundary established in Lemma~\ref{k:static-boundary}.
In the positive-lapse case its two commuting corner reflections
are carried to the coordinate reflections. Near
\(\vartheta=0\), their smooth even and odd parities give
\begin{equation}\label{k:corner-chart}
 \sigma=\lambda A_s(\lambda^2,y^2),\qquad
 \vartheta=yA_\vartheta(\lambda^2,y^2),\qquad
 A_s,A_\vartheta>0 ,
\end{equation}
with smooth coefficients; near the other pole use
\(\pi-\vartheta\) instead. This follows by dividing each odd
function by its coordinate and applying smooth even-function
factorization in the two variables. It does not require
analyticity. Thus \(\sigma/\lambda\), \(\sin\vartheta/y\), and
\(X/\sin\vartheta\) are smooth positive ratios jointly at the
corner. In the zero-lapse case the same polar conclusion follows
from the smooth chart in \((s,y)\), with the \(y\)-reflection.

For clarity, the end normalization in this construction retains an
initially unknown positive constant. We use the reflected
isothermal expansion \eqref{r:af-isothermal}, constructed in
Section~\ref{r:uniformization}. The estimate
\(Y=O_2(r^{-2})\) gives \(h_+-h=O_2(r^{-4})\) on the end.
After inversion of an asymptotically Euclidean meridional chart,
the reflected conformal structure extends across the origin with
H\"older coefficients of every exponent \(q'<1\).
Its isothermal rectification has a nonzero first derivative there;
rescaled annulus estimates give the differentiated remainder.
Reflection supplies the same expansion in ratio at the axes.
It follows that, for some \(c_\infty>0\),
\[
 \frac{X}{e^\sigma\sin\vartheta}\longrightarrow c_\infty,
 \qquad
 \frac{\sqrt{P_+}}{e^\sigma}\longrightarrow c_\infty
\]
uniformly, the first quotient interpreted by its continuous polar
extension. Since \(\lambda\to1\), the same first limit holds for
\(\rho\). The harmonicity proved in Lemma~\ref{k:timelike} becomes
flat harmonicity in the strip, and \(\rho\) vanishes continuously
on each finite side.

Put \(b'=2c_\infty\). For any \(\eps>0\), at a sufficiently distant
line \(\sigma=L\) the functions
\((b'\pm\eps)\sinh\sigma\sin\vartheta\) bound \(\rho\).
They have the same zero values on the finite sides. Maximum
comparison on the finite rectangles, followed by \(L\to\infty\)
and then \(\eps\downarrow0\), gives
\[
 \rho=b'\sinh\sigma\sin\vartheta .
\]
At the inner boundary, differentiation toward the domain gives
\(\partial_{\nu_+}\rho=\kappa X\), because
\(\rho=X\lambda\) and \(|\dd\lambda|_{h_+}=\kappa\).
As \(\nu_+=P_+^{-1/2}\partial_\sigma\), this says
\[
 X\sqrt{P_+}=\frac{b'}{\kappa}\sin\vartheta
                       \qquad(\sigma=0).
\]
The original and static metrics agree tangentially there.
The area of the original invariant boundary is therefore
\[
 A=2\pi\int_0^\pi X\sqrt{P_+}\,\dd\vartheta
     =\frac{4\pi b'}{\kappa}.
\]
Since \(\kappa=b/R_0^2\) and \(A=4\pi R_0^2\), we obtain \(b'=b\).
\end{proof}

\subsection{Identification of the map and the horizontal metric}

\begin{proposition}\label{k:identification}
In the coordinates of Lemma~\ref{k:rod-scale}, the stationary
map and horizontal metric coincide with those of the
\KN{} model of parameters \(M,a,Q\) from
Lemma~\ref{a:model-parameters}.
\end{proposition}

\begin{proof}
Let \(\Phi_*\) denote the model map in \eqref{a:model-map}, with
the same ordered axis constants as \(\Phi\), and set
\[
 w=\sinh\sigma\sin\vartheta,\qquad
 D=\dist_G(\Phi,\Phi_*)^2.
\]
Both maps satisfy the harmonic-map equation with weight \(w\).
The target is complete, simply connected, and nonpositively
curved, as established in the axial reduction. Joint convexity of
squared distance along its geodesics, applied to the two weighted
harmonic-map equations, gives
\begin{equation}\label{k:distance}
 \divg(w\nabla D)\ge0
\end{equation}
in the open strip; derivatives and divergence in this display are
Euclidean.

We verify the boundary controls needed to use this inequality.
By Lemma~\ref{a:axis-regularity}, on every bounded radial range
the two axial lengths are comparable
to the same ordinary transverse radius. Their logarithmic
difference is consequently bounded. If \(y\) is that radius, the
electromagnetic potential differences from their common axis
values are \(O(y^2)\), and the adjusted central potential
differences are \(O(y^4)\). In particular, writing
\(\Delta\psi=\psi-\psi_*\), \(\Delta\chi=\chi-\chi_*\), the central
displacement
\[
 (z-z_*)+\psi_*\Delta\chi-\chi_*\Delta\psi
\]
is \(O(y^4)\). To bound distance, first change the height
\(u=\log X\), then join the remaining three coordinates at fixed
height. The target length of the latter path is controlled by
\[
 X_*^{-1}(|\Delta\psi|+|\Delta\chi|)
 +X_*^{-2}|(z-z_*)+\psi_*\Delta\chi-\chi_*\Delta\psi|.
\]
These quantities are bounded, indeed vanish at a regular axis.
Lemma~\ref{k:static-boundary} preserves the ordinary transverse
radius by a smooth positive ratio at the horizon--axis corners,
so the same argument applies there. Away from the axes smoothness
of both maps gives boundedness directly. Thus \(D\) is bounded
on every compact radial range, including all finite boundary
faces and corners.

At infinity the matched scale in \eqref{k:rho} gives
\(u-u_*\to0\) uniformly. The end estimates \eqref{a:end-regularity} for the normalized
potentials, including their polar ratios, give respectively
\[
 X_*^{-1}(|\Delta\psi|+|\Delta\chi|)=O(e^{-\sigma}),
 \qquad
 X_*^{-2}|(z-z_*)+\psi_*\Delta\chi-\chi_*\Delta\psi|
                      =O(e^{-2\sigma}).
\]
Here the common axis constants remove the otherwise singular
polar terms. The preceding path estimate shows that \(D\to0\)
uniformly as \(\sigma\to\infty\).

Fix \(\delta>0\) and set \(v=(D-\delta)_+\).
It vanishes on a sufficiently distant tail and is bounded.
Test \eqref{k:distance} with \(\chi_{\mathrm{cut}}^2v\), where \(\chi_{\mathrm{cut}}\) is a
compactly supported cutoff in the open strip. The weak chain rule
and Cauchy's inequality give
\begin{equation}\label{k:capacity-test}
 \int w\chi_{\mathrm{cut}}^2|\nabla v|^2
       \le4\int wv^2|\nabla\chi_{\mathrm{cut}}|^2 .
\end{equation}
There is no outer error if the outer transition of \(\chi_{\mathrm{cut}}\) is
placed in the zero tail of \(v\). At each finite side, with
Euclidean distance \(x\) to that side, use a cutoff that rises
from zero to one on \(\eps^2<x<\eps\) with derivative
\((x|\log\eps|)^{-1}\). On a fixed radial range \(w\le Cx\), and
\[
 \int_{\eps^2}^{\eps}
       x\left|\frac1{x\log\eps}\right|^2\dd x
          =\frac1{|\log\eps|}\longrightarrow0.
\]
Products of these cutoffs treat the corners; their gradient
energies are bounded by the sum of the individual energies.
Letting \(\eps\downarrow0\) in \eqref{k:capacity-test} proves
\(\nabla v=0\) throughout the open strip. Its zero tail makes
\(v=0\). Finally \(\delta\downarrow0\) gives
\(\Phi=\Phi_*\).

It remains to determine a metric coefficient; map equality alone
does not yet do this. Put
\[
 q=\log(X\sqrt{P_+}),\qquad q_*=\log(X_*\sqrt P),
\]
where \(P=r^2+a^2\cos^2\vartheta\) is the model coefficient.
The static quotient has spatial metric
\(e^{2q}(\dd\sigma^2+\dd\vartheta^2)\). Its scalar constraint,
or directly \eqref{a:wave}, gives
\(-\Delta q=|\partial\Phi|_G^2\); hence \(q-q_*\) is harmonic.
It vanishes on \(\sigma=0\) by \eqref{k:area-density}.
At an ordinary axis the no-cone condition says
\(\sqrt{P_+}=|\partial_\vartheta X|\). Since \(X=X_*\),
the two nonsingular conformal scales have the same axis values.
Inside the strip their logarithmic difference is precisely
\[
 q-q_*=\frac12\log\frac{P_+}{P}.
\]
It extends boundedly through the corners, because the static
base metrics are smooth and positive, and tends uniformly to
zero at infinity: both square-root scales are asymptotic to
\((b/2)e^\sigma\). Maximum comparison now gives \(q=q_*\).
Thus \(P_+=P\), \(X=X_*\), and
\(\lambda=\rho/X=\lambda_*\), proving the proposition.
The oriented crossing order of the axes, future static time,
and increasing original azimuth have been preserved throughout;
the Hodge-star conventions in \eqref{a:quotient} therefore agree.
\end{proof}

\subsection{Reconstruction of the original data}

Return now to the original, unpolarized data
\((g,K,E_f,B_f)\). Its metric decomposition is
\[
 g=h+X^2(\dd\phi+\mathcal A)^2.
\]
For this subsection \(E',B'\) denote its constant duality rotation
with charges \((Q,0)\). The model's angular connection is
\(-\omega_*\,\dd t\), where
\[
 \omega_*=\frac{a(2Mr-Q^2)}{W},\qquad
 r=M+b\cosh\sigma,\qquad
 W=(r^2+a^2)^2-a^2\Delta\sin^2\vartheta .
\]

\begin{proposition}\label{k:lift}
The map of the regular interior orbits into the model given by
\begin{equation}\label{k:embedding}
 t=f,\qquad
 (\sigma,\vartheta)=(\sigma,\vartheta)(x),\qquad
 \phi_{\mathrm{KN}}=\phi+\beta,\qquad
 \dd\beta=\mathcal A+\omega_*\,\dd f
\end{equation}
is well defined and induces all of the original data, after inverse
constant duality rotation. It extends smoothly across the ordinary
interior axes and approaches spatial infinity.
\end{proposition}

\begin{proof}
Proposition~\ref{k:identification} identifies the horizontal
Lorentz metric, the map, and the map's normal velocities on the
graph \(t=f\). In particular, the curvature identity in
\eqref{a:quotient}, including
\[
 X^{-2}\omega(n)=s_3=-\frac{Xf_a}{2},
       \qquad \dd\mathcal A=f_a\,\dd A_h,
\]
identifies the pullback of \(\dd(-\omega_*\,\dd t)\) with
\(\dd\mathcal A\). Consequently
\(\dd(\mathcal A+\omega_*\,\dd f)=0\).
The meridian is simply connected, so the real primitive \(\beta\)
in \eqref{k:embedding} exists globally.
On the graph the model's fiber form becomes
\[
 \dd\phi_{\mathrm{KN}}-\omega_*\,\dd t
   =\dd\phi+\dd\beta-\omega_*\,\dd f
   =\dd\phi+\mathcal A.
\]
Its horizontal metric is
\(h_+-\lambda^2\dd f^2=h\) by \eqref{k:static}.
The induced metric is therefore exactly the original \(g\).

We also identify the full second fundamental form.
The identities \eqref{a:reconstruction} make this explicit.
The horizontal quotient tensor, together with \(n(\log X)\),
determines the horizontal block \(K|_h\), by the tensor relation
used in Proposition~\ref{j:adjoint}. The fiber-fiber unit entry is
\(K(e,e)=n(\log X)\). Finally the spatial identity
\[
 \omega=X^2*_h l,\qquad l=K(e,\cdot)|_h,
\]
determines the mixed block. All these quantities agree under the
map, with the same future normal, so no component of \(K\) is
discarded in undoing polarization.

Likewise, \(\dd\psi=X*_hE'_H\) and
\(\dd\chi=X*_hB'_H\) recover both horizontal electromagnetic
components. Their normal map velocities recover the remaining
components through \(E'_e=s_2\), \(B'_e=-s_1\).
Thus the rotated Maxwell fields agree. If \(Q>0\), apply
\begin{equation}\label{k:inverse-duality}
 E_f=\frac{Q_eE'-Q_bB'}Q,\qquad
 B_f=\frac{Q_bE'+Q_eB'}Q;
\end{equation}
if \(Q=0\), use the identity rotation. The model then has the
ordered original charges \((Q_e,Q_b)\).
The constant rotation preserves the stress tensor and
\(E_f\times B_f\), so it also preserves the total angular-momentum
normalization. The model choice \(a=-J/M\) consequently has exactly
the prescribed \(J\) with the positive convention for \(K\).

Every interior meridional change above respects axis reflection.
The differential defining the angular rotation has ordinary
invariant polar parity, so its primitive gives a smooth rotation
of each normal polar disk. This extends the map through the axes.
Finally \(\dd f=-\alpha_Y/\mu_*=O(r^{-2})\), and the global static
meridional coordinates tend to the model's spatial infinity.
\end{proof}

\subsection{Smooth attachment to the future horizon}

\begin{proposition}\label{k:horizon}
The map in Proposition~\ref{k:lift} extends to a smooth spacelike
embedding of the original exterior including \(S\), with a smooth
future unit normal. If \(N|_S=0\), its boundary is the bifurcation
sphere. If \(N|_S>0\), its boundary is a smooth cross-section of the
future horizon.
\end{proposition}

\begin{proof}
In the zero-lapse case the static meridional coordinates and
\(f\) are already smooth up to \(S\).
In the positive-lapse case, the isometric reflection in
Lemma~\ref{k:static-boundary}, followed by conformal reflection
across the inner circle, shows that \(\sigma\) is odd and
\(\vartheta\) even as functions of \((\lambda,z)\).
Hence
\[
 \frac{\sigma}{\lambda}>0,\qquad \vartheta
\]
extend as smooth even functions, and therefore as smooth
one-sided functions of the original coordinate \(\mu_*=\lambda^2\).
At the poles, the joint smooth positive ratios in
\eqref{k:corner-chart} give the original polar extension.

The apparent singularity of the angular variable is removed in
co-rotating coordinates. Put
\[
 \omega_H=\frac a{R_0^2},\qquad
 \phi_\sharp=\phi_{\mathrm{KN}}-\omega_H t.
\]
On the graph its angular adjustment has differential
\[
 \dd\beta-\omega_H\dd f
       =\mathcal A+(\omega_*-\omega_H)\dd f.
\]
The model expansion is
\(\omega_*-\omega_H=\sigma^2\omega_2(\sigma^2,\vartheta)\),
with smooth polar coefficients. This differential is therefore
smooth also in the positive-lapse case, where
\(\dd f=-\dd\mu_*/(2\kappa\mu_*)+\dd f_0\).
It is closed and has a smooth primitive up to the boundary,
with the same axis parity as before.

Use local horizon coordinates
\begin{equation}\label{k:null-coordinates}
 U=-\sigma e^{-\kappa t},\qquad
 V=\sigma e^{\kappa t}.
\end{equation}
They satisfy the exact identities
\begin{equation}\label{k:null-identities}
 -UV=\sigma^2,\qquad
 -\dd U\,\dd V=\dd\sigma^2-\kappa^2\sigma^2\dd t^2,\qquad
 \sigma^2\dd t=\frac{V\dd U-U\dd V}{2\kappa}.
\end{equation}
To check that these are smooth nonsingular extension coordinates,
the model has
\[
 \lambda_*^2=P\kappa^2\sigma^2
                       +\sigma^4 L(\sigma^2,\vartheta),
\]
where \(L\) and all the remaining model coefficients are smooth
in \(\sigma^2\) and the ordinary angular variables.
The leading horizontal block is
\[
 P(\dd\sigma^2-\kappa^2\sigma^2\dd t^2)
                  =-P\,\dd U\,\dd V.
\]
The error \(\sigma^4L\dd t^2\) is a smooth quadratic expression
in \(V\dd U-U\dd V\). The fiber form in co-rotating coordinates is
\(\dd\phi_\sharp-(\omega_*-\omega_H)\dd t\), which is smooth by
\eqref{k:null-identities}. Thus the complete Lorentz metric
extends nonsingularly, using ordinary angular patches at the
poles. For the electric model potential, its coefficient on the
horizon generator is constant. Subtract that constant times
\(\dd t\); the remaining time coefficient is \(O(\sigma^2)\),
smooth in \(\sigma^2\), so its field also extends smoothly.
The same holds after the constant duality rotation.

On the graph, if \(N|_S=0\), both \(U\) and \(V\) vanish smoothly
at \(S\). This is the bifurcation sphere. If \(N|_S>0\), use
\eqref{k:log-time} to write instead
\begin{equation}\label{k:future-graph}
 V=\frac{\sigma}{\sqrt{\mu_*}}e^{\kappa f_0},\qquad
 U=-\frac{\sigma}{\sqrt{\mu_*}}\,\mu_*e^{-\kappa f_0}.
\end{equation}
Both are smooth in the original variables, with
\(U=0\), \(V>0\) on \(S\). This is the future horizon branch:
in the exterior \(U<0<V\), and the future stationary generator
has components \(\kappa(-U\partial_U+V\partial_V)\).
Figure~\ref{k:horizon-figure} summarizes these two attachments.

The extended map preserves the smooth positive original metric.
It is thus an immersion and is spacelike up to \(S\);
the future unit normal then extends smoothly as well.
It is injective in the interior, since its static meridional
coordinates form a global coordinate change and each fiber map
is a rotation. On the boundary the map to the horizon
cross-section coordinates is a diffeomorphism, so distinct
boundary points remain distinct. No interior point maps to
the horizon. Finally compactness on bounded radial sets and
properness at the end make the injective map proper onto its
image, hence an embedding. The induced \(K,E_f,B_f\), already
identified in the interior, agree at \(S\) by smooth extension.
\end{proof}

\begin{figure}[htbp]
 \centering
 \begin{tikzpicture}[scale=0.95,>=Latex]
  \fill[blue!3] (-4,0) rectangle (0,3.7);
  \draw[->] (-4.3,0)--(0.8,0) node[below] {\(U\)};
  \draw[->,thick] (0,-0.35)--(0,4.0) node[above] {\(V\)};
  \draw[very thick,blue!65!black]
      (-3.25,3.35) .. controls (-1.5,2.15) and (-0.55,1.65)
                   .. (0,1.45);
  \draw[thick,dashed,teal!75!black] (-3.1,3.1)--(0,0);
  \fill[blue!65!black] (0,1.45) circle (2pt);
  \fill (0,0) circle (2pt);
  \node[align=left,anchor=west] (future) at (0.55,2.25)
      {\(N|_S>0\)\\\(U=0,\ V>0\)};
  \draw[->] (future.west)--(0.08,1.49);
  \node[align=left,anchor=west] (bif) at (0.75,0.65)
      {\(N|_S=0\)\\\(U=V=0\)};
  \draw[->] (bif.west)--(0.08,0.03);
  \node[align=center] at (-2.45,0.75)
      {exterior\\\(U<0<V\)};
 \end{tikzpicture}
 \caption{The two smooth horizon attachments in
 \eqref{k:null-coordinates}, with angular directions suppressed.
 The curves depict alternative spacelike slices. In the
 positive-lapse case the boundary value of \(V\) may vary over
 the horizon sphere, as in \eqref{k:future-graph}.}
 \label{k:horizon-figure}
\end{figure}

\subsection{The converse for arbitrary admissible slices}

For a \KN{} spacetime whose parameters have already been
identified with \(m,J,Q_e,Q_b\), every smooth cross-section of
the indicated horizon has area
\begin{equation}\label{k:horizon-area}
 A_H=4\pi\left[
 \left(m+\sqrt{m^2-(J/m)^2-Q^2}\right)^2+(J/m)^2\right].
\end{equation}
Indeed the horizon is nonexpanding and its null induced metric
descends to its sphere of generators. The area is independent
of the section. Substitution in the parameter identity of
Lemma~\ref{a:model-parameters} gives equality in Theorem~\ref{m:main}.
This applies to nonconstant time graphs as well.

We supply the end identification also when an embedding into a
subextremal \KN{} exterior is given without initially equating
its ambient parameters with the fluxes of the slice. This
ensures that the converse uses exactly the asymptotic frame in
the theorem.

\begin{lemma}\label{k:end-identification}
Let an original, unpolarized data set satisfying the asymptotic
and rest-frame hypotheses of Theorem~\ref{m:main} be realized
as a spacelike slice approaching spatial infinity in a
subextremal \KN{} exterior with parameters
\(M,\widetilde a,\widetilde Q_e,\widetilde Q_b\). Then
\[
 m=M,\qquad J^2=M^2\widetilde a^2,\qquad
 Q_e^2+Q_b^2=\widetilde Q_e^2+\widetilde Q_b^2 .
\]
In a sufficiently distant stationary spatial chart the slice
is a single graph \(t=F(\bar x)\), with
\(\dd F=O_1(|\bar x|^{-2})\).
\end{lemma}

\begin{proof}
Write \(r_x=|x|\) in the prescribed slice chart, let \(T\) be
the stationary time Killing field future at infinity, and
decompose
\[
 T=N_Tn+Y_T .
\]
Since the end approaches ambient spatial infinity,
\(N_T>|Y_T|_g\) eventually and
\(N_T^2-|Y_T|_g^2=-\bar g(T,T)\to1\).
We first obtain
\begin{equation}\label{k:restricted-killing}
 N_T\to1,\qquad Y_T=O_2(r_x^{-2}).
\end{equation}
No prior bound on the tilt of the embedding is needed.
The Gauss equation bounds the tangential ambient curvature
components in the slice frame by \(O(r_x^{-3})\); Codazzi
bounds the components with one normal index by the same order.
The remaining normal-normal block follows by contraction
with the ambient Ricci tensor, whose components in this
frame are \(O(r_x^{-4})\) by the given electromagnetic fields
and the Einstein--Maxwell equation. Thus
\(\overline{\operatorname{Rm}}=O(r_x^{-3})\) in the entire
slice frame.

The Killing first-jet equation
\(\bar\nabla^2T=\overline{\operatorname{Rm}}*T\), split into
tangential and normal components, therefore has the end
bounds used in Lemma~\ref{j:end}. The splitting adds only
\(K*\bar\nabla T\) and
\((\nabla K+K*K)*T\); the skew symmetry of \(\bar\nabla T\)
controls all its components by the tangential derivatives.
In slice coordinates the Christoffel symbols and their
derivatives are \(O(r_x^{-2})\), \(O(r_x^{-3})\).
The radial first-jet integration and causality argument of
that lemma consequently give constant limits with
\(O_2(r_x^{-q'})\) errors, for any fixed \(1/2<q'<1\).
Here the normalization is determined by
\(-\bar g(T,T)\to1\), rather than by a variational prototype.

For the original unpolarized data, the constraints and prescribed falloff give
\[
 L(g-\delta)=\partial_i\partial_j(g-\delta)_{ij}
       -\Delta_\delta\tr_\delta(g-\delta)=O(r_x^{-4}).
\]
The pair \((N_T,Y_T)\) satisfies
\(\operatorname{sym}\nabla Y_T=-N_TK\), with \(K=O_1(r_x^{-3})\),
bounded \(N_T\), and the constant asymptotics just proved. Since \(m>0\),
the general assertion of Lemma~\ref{j:end} gives \(Y_{T,\infty}=0\).
The norm limit then gives \(N_{T,\infty}=1\).
Prolonging \(\operatorname{sym}\nabla Y_T=-N_TK\) gives
\[
 |\partial^2Y_T|
 \le C r_x^{-2}|\partial Y_T|
       +Cr_x^{-3}|Y_T|+Cr_x^{-4}.
\]
Twice integrating from the zero end limits, and reinserting
the improved bounds, yields \eqref{k:restricted-killing}.

Project along \(T\) to stationary spatial coordinates
\(\bar x\) of the ambient exterior, and restrict to
\(\bar r=|\bar x|\) sufficiently large. This projection is a
local diffeomorphism on the slice, because a timelike vector
cannot be tangent to a spacelike hypersurface. It is proper
over that spatial tail: the end goes to ambient spatial
infinity, and a compact core has bounded projected radius.
Explicitly, for every ambient radius \(R\) there is a slice
radius \(R_x(R)\) such that
\[
 |x|>R_x(R)\quad\Longrightarrow\quad |\bar x|>R.
\]
Thus the preimage of the ambient tail contains one whole
connected intrinsic coordinate tail, and the preimage of
any compact subset of the ambient tail is a closed subset
of an intrinsic compact set.
Its image is therefore open and closed in the connected
tail, so it is surjective and a covering. The spatial tail
is simply connected. Each component of its preimage covers
the whole tail and hence reaches arbitrarily far out in
the slice end: otherwise its closure would lie in an
intrinsic compact set, whose projected radius is bounded.
All such components must therefore intersect the one
connected full intrinsic tail just specified. There is just
one component and one sheet. This proves the single-graph
assertion.

The stationary orbit metric on this three-dimensional end
pulls back to
\begin{equation}\label{k:orbit-metric}
 g_T=g+\frac{(Y_T^\flat)^2}{N_T^2-|Y_T|^2}.
\end{equation}
By \eqref{k:restricted-killing}, it differs from \(g\) by
\(O_2(r_x^{-4})\) and has the same mass.
In the ambient stationary chart, the same metric differs
from the canonical spatial metric by
\(\bar g_{ti}\bar g_{tj}/(-\bar g_{tt})=O_2(\bar r^{-4})\).
The canonical mass is \(M\).

Here the equality of masses in the two charts can be checked
directly. Since both represent the same uniformly Euclidean
metric, \(D\bar x\) and its inverse are bounded. Distances in
the two ends show that \(r_x\) and \(\bar r\) are comparable.
The transformation law for the Christoffel symbols then
gives \(D^2\bar x=O(r_x^{-2})\). Integrating along rays and
sphere paths yields a constant orthogonal matrix \(O_*\) with
\begin{equation}\label{k:chart-transition}
 D\bar x=O_*+O(r_x^{-1}),\qquad
 \bar x=O_*x+O(\log r_x).
\end{equation}
The ADM flux comparison separates a pure Euclidean
symmetric gradient, whose linear flux vanishes, from the
original perturbation. Quadratic errors tend to zero.
The Euclidean scalar linearization is integrable
(\(O(r_x^{-4})\)), so asymptotically round deformed spheres
give the same limit as round spheres. Thus \(m=M\).

In the \(\bar x\)-chart the graph satisfies
\begin{equation}\label{k:graph-gradient}
 \dd F=
 \frac{\bar g_{ti}\dd\bar x^i-Y_T^\flat}
                  {N_T^2-|Y_T|^2}
             =O_1(\bar r^{-2}).
\end{equation}
Its vanishing tilt leaves the leading Coulomb fluxes
unchanged, giving
\(Q_e^2+Q_b^2=\widetilde Q_e^2+\widetilde Q_b^2\);
possible orientation changes do not affect this identity.
The graph tangents and the normalized future normal
covector proportional to \(-\dd t+\dd F\) give, with the
positive convention for \(K\),
\begin{equation}\label{k:graph-tensor}
 K_{ij}=K^{\rm can}_{ij}+\partial_i\partial_jF
                                  +O(\bar r^{-4}).
\end{equation}
For every Euclidean rotation vector \(R\),
\[
 n_\delta^i(\partial_i\partial_jF)R^j
       =R\left(\partial_{\bar r}F-\frac{F}{\bar r}\right).
\]
Its round-sphere integral is zero because \(R\) preserves
the sphere area form. The trace-reversed Euclidean term
also vanishes, since \(R\) is tangential; all metric,
normal, and area errors tend to zero at these decay orders.
Hence the angular flux vector in this chart has the
canonical norm \(|M\widetilde a|\).

Under \eqref{k:chart-transition}, rotation vectors agree,
up to the constant orthogonal change, with errors
\(O(\log r_x)\). Their flux errors are
\(O(\log r_x/r_x)\), since \(K=O_1(r_x^{-3})\).
The corresponding spheres differ in radius by
\(O(\log r_x)\), and their normals and relative area
elements by \(O(\log r_x/r_x)\); these also give vanishing
angular-flux errors. In the prescribed slice chart,
axisymmetry forces the two transverse components of the
angular flux vector to vanish. Its axial component is
the prescribed total \(J\), because the electromagnetic
correction tends to zero at infinity. Thus its norm is
\(|J|\), proving \(J^2=M^2\widetilde a^2\).
\end{proof}

If an admissible embedded slice is only said to meet the
appropriate future horizon, its boundary is nevertheless a
cross-section. An interior point cannot lie on the null face,
by spacelikeness and the fact that the whole image lies on its
exterior side. At a boundary contact use the smooth horizon
coordinate \(U\), with \(U=0\) on the future branch,
\(U<0\) in the exterior, and \(\dd U\) positive on future
timelike vectors. Extend the spacelike embedding locally a
little beyond its boundary; spacelikeness is an open condition.
The cut \(U=0\) of this extension is a smooth MOTS with normal
toward \(U<0\): its outgoing future null direction is tangent
to the nonexpanding horizon. This remains true at bifurcation.
The given \(S\) touches this cut from the exterior side with
the same outward orientation. Local graph strong comparison
for the expansion equation makes the surfaces coincide near
the contact. The contact set is therefore open and closed
in the connected \(S\), so all of \(S\) lies on that branch.
Projection along the horizon generators is a local
diffeomorphism by spacelikeness. Compactness makes it a
covering of the generator sphere, and simple connectivity
makes it a diffeomorphism. Formula~\eqref{k:horizon-area},
with the parameter identifications of
Lemma~\ref{k:end-identification}, now proves the converse.

\begin{lemma}[Exclusion of partial bifurcation contact]
\label{k:horizon-cut-dichotomy}
A smooth spacelike cross-section of the closed future horizon bounding an exterior
with \(\theta_+=0\) and the all-cut outer area-minimizing property
lies entirely on the open future horizon or is the bifurcation
sphere.
\end{lemma}

\begin{proof}
This is a direct boundary argument, independent of the adjoint
boundary dichotomy. In the smooth co-rotating coordinates
\eqref{k:null-coordinates}, write the section as
\(U=0,\ V=v(x)\ge0\), where \(x\) ranges over the sphere of
generators. Such a global graph exists because generator
projection is a diffeomorphism.
Smoothness of the model coefficients in \(-UV\), with the
weights imposed by the horizon Killing field
\(\kappa(-U\partial_U+V\partial_V)\), gives on \(U=0\)
\[
 \bar g_{UV}=-c,\quad
 \bar g_{UA}=Vb_A,\quad
 \bar g_{AB}=q_{AB},\quad
 \bar g_{VA}=\bar g_{VV}=0,\qquad c>0,
\]
and
\[
 \partial_U\bar g_{AB}=Vd_{AB},\qquad
 \partial_U\bar g_{VA}=e_A,\qquad
 \partial_U\bar g_{VV}=0.
\]
Here \(q\) is a positive sphere metric and all coefficients
are smooth functions of \(x\). The last identity follows from
the stronger form \(\bar g_{VV}=U^2a_2(UV,x)\); merely knowing
\(\bar g_{VV}=0\) on the horizon would not suffice.

Set \(\ell=\partial_V\), and choose the transverse future null
normal \(k\) with \(\bar g(k,\ell)=-1\). The section tangents
are \(e_A^{\,S}=\partial_A+v_A\partial_V\); orthogonality gives
\[
 k^U=\frac1c,\qquad
 k^A=q^{AB}\left(v_B-\frac{v b_B}{c}\right),\qquad
 k_V=-1,\qquad k_A=v_A.
\]
The remaining component is fixed by nullness and is immaterial
for the expansion because the horizon's outgoing second
fundamental form is zero. Direct contraction of
\(-\bar g(k,\bar\nabla_{e_A^{\,S}}e_B^{\,S})\) with \(q^{AB}\)
gives the linear expression
\begin{equation}\label{k:cut-expansion}
 \theta_k[v]=
 \Delta_q v+
 \frac1c\langle\dd c+e-b,\dd v\rangle_q+
 \frac{\tr_qd-2\operatorname{div}_q b}{2c}\,v .
\end{equation}
For example the relevant lower Christoffel coefficients are
\[
 \Gamma_{U;AB}
    =\frac V2(\partial_A b_B+\partial_B b_A-d_{AB}),\qquad
 \Gamma_{U;VB}
    =\frac12(b_B-\partial_Bc-e_B);
\]
combining these with \(k_V=-1\), \(k_A=v_A\) yields
\eqref{k:cut-expansion}. The \(VV\) coefficients vanish on
the horizon, so there is no quadratic gradient term.

Outer area minimization, applied to arbitrary nonnegative
outward variations, implies \(H\ge0\). Since
\(\theta_+=H+\tr_SK=0\), the other future expansion is
\(\theta_-=\tr_SK-H=-2H\le0\).
The outward future normal \(n+\nu\) is a positive multiple
of \(\ell\), so \(n-\nu\) is a positive multiple of \(k\).
Consequently \(\theta_k[v]\le0\).
Equation~\eqref{k:cut-expansion} is a uniformly elliptic
linear inequality for the nonnegative smooth function \(v\).
The strong minimum principle applies even if its zeroth-order
coefficient has either sign: in the operator
\(-\theta_k\), replace that coefficient by its positive part,
which preserves the supersolution inequality because
\(v\ge0\). Hence a zero of \(v\) forces \(v\equiv0\) on the
connected sphere. Otherwise \(v>0\) everywhere, proving the
two alternatives.
\end{proof}

\subsection{Sharp examples and completion of the proof}

The canonical constant Boyer--Lindquist-time slice
\(\sigma\ge0\) provides a sharp example for every subextremal
choice of parameters on the strict branch. Its topology is
the rotation exterior, and it is complete with its compact
boundary included. In the Euclidean end radius
\(\rho_0=(b/2)e^\sigma\), the model metric has Cartesian
expansion
\[
 g^{\rm can}_{ij}
   =\left(1+\frac{2M}{\rho_0}\right)\delta_{ij}
                                        +O_4(\rho_0^{-2}).
\]
The angular shift is \(O(\rho_0^{-3})\), so
\(K^{\rm can}=O_3(\rho_0^{-3})\). The electric model potential
has time coefficient \(Q/r+O(r^{-3})\) and spatial angular
coefficient \(O(r^{-1})\sin^2\vartheta\). It gives the required
radial Coulomb electric term and \(O_3(\rho_0^{-3})\) magnetic
field; inverse duality gives both prescribed charges.
All differentiated bounds follow by smooth expansion in
\(\rho_0^{-1}\) with ordinary polar parity. The bifurcate
coordinates used above prove inner-boundary regularity.
The source-free constraints, integrability, and
\(P_{\rm ADM}=0\) follow, with energy \(M\), angular momentum
\(-aM\), and the specified charges.

We verify the two geometric boundary hypotheses for these
examples, including cuts that are neither invariant nor
trapped. On the \(\sigma\)-spheres, the tangential trace of
\(K^{\rm can}\) is zero, because its only nonzero entries
are mixed azimuthal entries. For the smooth outward unit
radial field \(V_r=P^{-1/2}\partial_\sigma\),
\begin{equation}\label{k:radial-expansion}
 \theta_+=H=\divg V_r
       =\frac{\partial_\sigma W}{2W\sqrt P}.
\end{equation}
This is zero at \(\sigma=0\) and positive outside. Indeed,
\[
 \partial_rW
 =4r(r^2+a^2)-2a^2(r-M)\sin^2\vartheta
 \ge4r^3+2a^2(r+M)>0 ,
\]
and \(\partial_\sigma r=b\sinh\sigma>0\) outside.
At a largest-radius point of an enclosing interior surface,
the outward tangent normal is \(V_r\) and its mean curvature
is at least that of the tangent radial leaf. The tangential
planes, hence their \(K\)-traces, agree at that point.
Equation~\eqref{k:radial-expansion} therefore excludes every
compact enclosing surface with \(\theta_+\le0\) everywhere,
also when it has several components.

Finally let \(D\) be any enclosing end-side domain in the
definition of Theorem~\ref{m:main}, with full intrinsic
boundary \(\Gamma\), and let \(n_D\) point out of \(D\).
Truncate the whole exterior and \(D\) at a common distant
radial sphere. Subtracting their divergence identities gives
\[
 0\le\int_{\Omega\setminus D}\divg V_r\,\dd V_g
      =-A(S)-\int_\Gamma\langle V_r,n_D\rangle\,\dd A_g.
\]
The identity counts portions of \(\Gamma\) coinciding with
\(S\), including the case \(D=\Omega\). Since \(|V_r|=1\),
\[
 A(S)\le-\int_\Gamma\langle V_r,n_D\rangle\,\dd A_g
                                         \le\Area_g(\Gamma).
\]
Thus the canonical horizon is outer area-minimizing for the
full class of cuts.

\begin{proof}[Completion of Theorem~\ref{m:main}]
Corollary~\ref{d:numerical} proves the numerical inequality on
the closed physical branch. On the strict branch,
Proposition~\ref{j:adjoint} supplies the stationary quotient
and its boundary conditions. Lemmas~\ref{k:timelike}--\ref{k:rod-scale}
and Proposition~\ref{k:identification} identify that quotient
with the subextremal model determined by
Lemma~\ref{a:model-parameters}. Propositions~\ref{k:lift}
and~\ref{k:horizon} give the required smooth embedding of
the original metric, tensor, and electromagnetic fields,
with its future normal and its stated boundary image.
The horizon area formula and
Lemma~\ref{k:end-identification} prove the converse, and the
canonical examples verify sharpness within the stated class.
No classification on the extremal equality branch is asserted.
\end{proof}
